% Part: first-order-logic
% Chapter: axiomatic-proofs
% Section: identity

\documentclass[../../../include/open-logic-section]{subfiles}

\begin{document}

\olfileid{fol}{axd}{ide}

\olsection{\usetoken{P}{derivation} mawa \usetoken{S}{identity}}

Kanggo nampung $\eq$ ing !!{derivation}s, kita mung perlu nambah
skema-skema aksioma anyar. Definisi !!{derivation} lan $\Proves$ tetep
padha; kita mung uga ngidini aksioma-aksioma anyar kasebut.

\begin{defn}[Aksioma kanggo !!{identity}]
\begin{align}
\ollabel{ax:id1} & \eq[t][t], \\
\ollabel{ax:id2} & \eq[t_1][t_2] \lif (!B(t_1) \lif
  !B(t_2)),
\end{align}
kanggo saben term katutup $t$, $t_1$, $t_2$.
\end{defn}

\begin{prop}\ollabel{prop:sound}
  Aksioma \olref{ax:id1} lan \olref{ax:id2} iku valid.
\end{prop}

\begin{proof}
  Gladhen.
\end{proof}

\begin{prob}
Buktekna \olref[fol][axd][ide]{prop:sound}.
\end{prob}

\begin{prop}\ollabel{prop:iden1}
 $\Gamma \Proves \eq[t][t]$, kanggo saben term katutup $t$ lan
 himpunan~$\Gamma$.
 Cathetan koreksi sumber OLPL-040: watesan term tertutup diwarisake
 saka skema aksioma identitas kapisan.
\end{prop}

\begin{prop}\ollabel{prop:iden2}
  Kanggo term-term katutup, yen $\Gamma \Proves !A(t_1)$ lan $\Gamma
  \Proves \eq[t_1][t_2]$, mula $\Gamma \Proves !A(t_2)$. Cathetan
  koreksi sumber OLPL-040: watesan iki dibutuhake amarga bukti
  nggunakake skema aksioma identitas kapindho.
\end{prop}

\begin{proof}
!!{formula}
\[
(\eq[t_1][t_2] \lif (!A(t_1) \lif !A(t_2)))
\]
iku conto saka~\olref{ax:id2}. Kesimpulan kasebut ndherek saka rong
aplikasi~\MP.
\end{proof}

\end{document}
